\documentclass[final]{beamer}

\usepackage[size=custom,width=100,height=180,scale=2.0]{beamerposter}
\usepackage{amsmath,amssymb,amsthm}
\usepackage{tikz}
\usetikzlibrary{arrows.meta,positioning,calc,decorations.pathreplacing,shapes.geometric,fit,backgrounds}
\usepackage{booktabs}
\usepackage{multicol}
\usepackage{xcolor}
\usepackage{tcolorbox}
\tcbuselibrary{skins,breakable,raster}

\definecolor{mainblue}{RGB}{31,73,125}
\definecolor{accentblue}{RGB}{68,138,201}
\definecolor{lightblue}{RGB}{213,232,252}
\definecolor{accentred}{RGB}{192,57,43}
\definecolor{accentgreen}{RGB}{39,128,80}
\definecolor{accentpurple}{RGB}{102,51,153}
\definecolor{boxbg}{RGB}{245,248,255}
\definecolor{titlebg}{RGB}{20,50,100}
\definecolor{demandcol}{RGB}{210,120,0}

\usetheme{default}
\setbeamercolor{background canvas}{bg=white}
\setbeamercolor{headline}{bg=titlebg,fg=white}
\setbeamercolor{footline}{bg=titlebg,fg=white}
\setbeamerfont{headline}{size=\Large}

\tcbset{
  posterbox/.style={
    enhanced,
    colback=boxbg,
    colframe=mainblue,
    arc=4pt,
    boxrule=1.5pt,
    fonttitle=\bfseries\large,
    coltitle=white,
    attach boxed title to top left={yshift=-2mm,xshift=4mm},
    boxed title style={colback=mainblue,arc=3pt,boxrule=0pt,boxsep=1.1mm,top=1.2mm,bottom=1.2mm},
    top=6mm,
    left=4mm,right=4mm,bottom=4mm,
  },
  defbox/.style={
    enhanced,
    colback=lightblue,
    colframe=accentblue,
    arc=4pt,
    boxrule=1.5pt,
    fonttitle=\bfseries\normalsize,
    coltitle=white,
    attach boxed title to top left={yshift=-2mm,xshift=4mm},
    boxed title style={colback=accentblue,arc=3pt,boxrule=0pt,boxsep=1.1mm,top=1.2mm,bottom=1.2mm},
    top=6mm,
    left=4mm,right=4mm,bottom=4mm,
  },
  thmbox/.style={
    enhanced,
    colback=white,
    colframe=accentgreen,
    arc=4pt,
    boxrule=1.5pt,
    fonttitle=\bfseries\normalsize,
    coltitle=white,
    attach boxed title to top left={yshift=-2mm,xshift=4mm},
    boxed title style={colback=accentgreen,arc=3pt,boxrule=0pt,boxsep=1.1mm,top=1.2mm,bottom=1.2mm},
    top=6mm,
    left=4mm,right=4mm,bottom=4mm,
  },
  appbox/.style={
    enhanced,
    colback=white,
    colframe=accentpurple,
    arc=4pt,
    boxrule=1.5pt,
    fonttitle=\bfseries\normalsize,
    coltitle=white,
    attach boxed title to top left={yshift=-2mm,xshift=4mm},
    boxed title style={colback=accentpurple,arc=3pt,boxrule=0pt,boxsep=1.1mm,top=1.2mm,bottom=1.2mm},
    top=6mm,
    left=4mm,right=4mm,bottom=4mm,
  },
}

\newcommand{\br}[1]{\left(#1\right)}
\newcommand{\brc}[1]{\left\{#1\right\}}
\newcommand{\spr}[1]{\left|#1\right|}
\newcommand{\fl}[1]{\left\lfloor#1\right\rfloor}

\newcommand{\defterm}[2]{%
  \begin{tcolorbox}[defbox,title={Definition: #1}]
  #2
  \end{tcolorbox}}

\newcommand{\thmterm}[2]{%
  \begin{tcolorbox}[thmbox,title={#1}]
  #2
  \end{tcolorbox}}

\begin{document}
\begin{frame}[t]{}

\begin{tcolorbox}[
  enhanced,colback=titlebg,colframe=titlebg,arc=0pt,boxrule=0pt,
  left=12mm,right=12mm,top=12mm,bottom=12mm
]
  {\color{white}\bfseries\fontsize{174}{202}\selectfont
   Graph Partitioning with Demands:\par
   \fontsize{78}{92}\selectfont
   Generalized Graph Conductance with Applications to Hierarchical Clustering}
  \par
  {\color{accentblue!80!white}\large
   Micha{\l} Szyfelbein \quad \& \quad Dariusz Dereniowski\par
   Gda\'nsk University of Technology}
\end{tcolorbox}


\begin{columns}[T]

\begin{column}{0.48\textwidth}

\begin{tcolorbox}[posterbox,title={Motivation}]
\textbf{Sparsest cut} is to find a cut $\br{S, \bar{S}}$, minimizing
\[
  \phi(S) = \frac{c(S,\bar S)}{|S|\cdot|\bar S|}.
\]
This problem is important due to its usage in graph partitioning,
divide-and-conquer algorithms, and hierarchical clustering.

\medskip
The denominator $|S|\cdot|\bar S|$ has \emph{two} interpretations:
\begin{itemize}
  \item \textbf{Demand view}: total unit demand crossing the cut.
  \item \textbf{Balance view}: large when both sides are of similar size.
\end{itemize}

\medskip
With a general demand function $w\colon V\times V\to\mathbb{N}$ these two
interpretations \textbf{diverge}.  Existing work
only regards minimising $\phi_w(S)=c(S,\bar S)/w(S,\bar S)$, \emph{ignores the balance of the cut}.
\end{tcolorbox}



\defterm{Graph with demands}{
  A graph $G=(V,E,c,w)$ where $c\colon E\to\mathbb{N}$ is a capacity
  (edge-cut cost) function and $w\colon V\times V\to\mathbb{N}$ is a
  \emph{demand} function.  For disjoint $A,B\subseteq V$:
  \[
    w(A,B)=\sum_{u\in A}\sum_{v\in B}w(u,v),\qquad
    w(A,V)=w(A)+w(A,\bar A).
  \]
}



\begin{tcolorbox}[posterbox,title={Illustration 1: Weighted Degree of $S$},top=12mm,bottom=12mm]
\centering
\resizebox{0.98\linewidth}{!}{%
\begin{tikzpicture}[>=Stealth,line width=1.4pt,
  vn/.style={circle,draw=mainblue,fill=white,minimum size=5mm,inner sep=0pt,line width=1pt,
         preaction={fill=black,opacity=0.18,transform canvas={xshift=0.8pt,yshift=-0.8pt}}}]

  \draw[fill=accentblue!25,draw=accentblue,line width=2pt]
    (0,0) ellipse (3.2 and 1.8);
  \node[accentblue,font=\bfseries\normalsize,draw=none,fill=none] at (0,2.8) {$S$};
  \node[vn] (s1) at (-1.4, 0.5) {};
  \node[vn] (s2) at ( 0.0, 0.9) {};
  \node[vn] (s3) at ( 1.2, 0.3) {};
  \node[vn] (s4) at (-0.3,-0.8) {};
  \draw[accentblue!70] (s1)--(s2)--(s3) (s1)--(s4)--(s3) (s2)--(s4);

  \draw[fill=accentgreen!25,draw=accentgreen,line width=2pt]
    (8,0) ellipse (3.2 and 1.8);
  \node[accentgreen,font=\bfseries\normalsize,draw=none,fill=none] at (8,2.8) {$\bar S$};
  \node[vn] (b1) at ( 6.6, 0.6) {};
  \node[vn] (b2) at ( 8.0, 1.0) {};
  \node[vn] (b3) at ( 9.2, 0.4) {};
  \node[vn] (b4) at ( 7.8,-0.8) {};
  \draw[accentgreen!70] (b1)--(b2)--(b3) (b2)--(b4)--(b3);

  \draw[demandcol,line width=1.4pt,dashed] (s3)--(b1);
  \draw[demandcol,line width=1.4pt,dashed] (s2)--(b2);
  \draw[demandcol,line width=1.4pt,dashed] (s3)--(b4);
  \draw[demandcol,line width=1.4pt,dashed] (s4)--(b1);

  \node[draw=none,fill=none,font=\normalsize,anchor=west,align=left] at (12.2,0.4) {
    $w(S,V) =$
    $\textcolor{accentblue}{w(S)}$
    $+ \textcolor{demandcol}{w(S,\bar S)}$
  };

\end{tikzpicture}}
\end{tcolorbox}




\defterm{Generalized Sparsest Cut (Arora-Lee-Naor)}{
  \[
    \phi_w(S) = \frac{c(S,\bar S)}{w(S,\bar S)}.
  \]
  Best known approximation ratio: $O\!\br{\sqrt{\log n}\cdot\log\log n}$.
}



\defterm{Generalized Conductance \textnormal{[this work]}}{
  \[
    \psi_w(S) = \frac{c(S,\bar S)}{w(S,V)\cdot w(\bar S,V)}.
  \]
  The denominator uses \emph{ weighted degree}, measuring
  balance even when little demand crosses the cut.
}



% --- Figure 1: Auxiliary graph G' for Step 1 ---
\begin{tcolorbox}[posterbox,title={Illustration 2: Auxiliary Graph $G'$ (Step~1)}]
\small
A root $r$ is added to $G$ with zero-capacity edges, thus forming $G'$.
Its demand to each $u\in V$ equals $w'(r,u)=w(u,V)$,
encoding the weighted degree of $u$. This then becomes an instance of the generalized sparsest-cut, under an \textbf{additional condition}, that $w'\br{S, \bar{S}}\leq K$, where we set $K=\frac{4}{3}w\br{V}$. 

\vspace{3mm}
\centering
\resizebox{0.95\linewidth}{!}{%
\begin{tikzpicture}[>=Stealth, line width=1.0pt,
  vn/.style={circle,draw=mainblue,fill=white,
         minimum size=9mm,inner sep=0pt,line width=1.0pt,font=\tiny,
         preaction={fill=black,opacity=0.18,transform canvas={xshift=0.8pt,yshift=-0.8pt}}}]

  \node[vn,draw=accentred,fill=accentred!20,
        minimum size=10mm,line width=1.5pt] (r) at (-11.0, 0) {$r$};
  \node[accentred,font=\tiny,draw=none,fill=none,
        align=center,anchor=north] at (-11.0,-1.1)
    {$c(r,{\cdot})=0$};

  \draw[fill=lightblue!55,draw=accentblue,line width=1.5pt]
    (6.0,0) ellipse (6.5 and 5.0);
  \node[accentblue,font=\bfseries\tiny,draw=none,fill=none] at (6.0,5.7)
    {$G$};

  \node[vn] (u) at (4.0,  3.5) {$u$};
  \node[vn] (v) at (8.0,  0.0) {$v$};
  \node[vn] (x) at (4.0, -3.5) {$x$};
  \draw[mainblue!70,line width=1.0pt] (u)--(v)--(x)--(u);

  \draw[dashed,accentred,line width=1.2pt] (r) -- (u)
    node[midway,above,sloped,draw=none,fill=none,font=\tiny,accentred]
      {$w'(r,u){=}w(u,V)$};
  \draw[dashed,accentred,line width=1.2pt] (r) -- (v)
    node[midway,above,sloped,draw=none,fill=none,font=\tiny,accentred]
      {$w'(r,v){=}w(v,V)$};
  \draw[dashed,accentred,line width=1.2pt] (r) -- (x)
    node[midway,below,sloped,draw=none,fill=none,font=\tiny,accentred]
      {$w'(r,x){=}w(x,V)$};

\end{tikzpicture}}


All other demand values (inbetween original vertices) become $0$.
\vspace{5mm}
\end{tcolorbox}




% --- Figure 3: Cut-sparsifier tree decomposition ---
\begin{tcolorbox}[posterbox,title={Illustration 3: Cut-Sparsifier Decomposition}]
\small
To solve the constrained sparsest-cut, embed $G'$ into a convex combination of trees, such that the expected value of each cut is preserved up to an $O(\log n)$ distortion.
Then, solve the constrained sparsest cut \emph{exactly} on each tree, move each cut to $G'$ and output the best solution.

\vspace{3mm}
\centering
\resizebox{0.90\linewidth}{!}{%
\begin{tikzpicture}[>=Stealth, line width=1.2pt,
  vn/.style={circle,draw=mainblue,fill=white,
         minimum size=5mm,inner sep=0pt,line width=1.2pt,
         preaction={fill=black,opacity=0.18,transform canvas={xshift=0.8pt,yshift=-0.8pt}}},
  tn/.style={circle,draw=accentgreen,fill=accentgreen!18,
         minimum size=4mm,inner sep=0pt,line width=1.2pt,
         preaction={fill=black,opacity=0.18,transform canvas={xshift=0.8pt,yshift=-0.8pt}}}]

  \node[vn] (g1) at ( 0.0, 1.5) {};
  \node[vn] (g2) at ( 1.5, 3.0) {};
  \node[vn] (g3) at ( 3.0, 2.5) {};
  \node[vn] (g4) at ( 4.5, 1.0) {};
  \node[vn] (g5) at ( 3.5,-1.0) {};
  \node[vn] (g6) at ( 1.5,-1.5) {};
  \node[vn] (g7) at (-0.5,-0.5) {};
  \node[vn] (g8) at ( 2.2, 0.8) {};
  \draw[mainblue!75,line width=1.6pt]
    (g1)--(g2)--(g3)--(g4)--(g5)--(g6)--(g7)--(g1)
    (g2)--(g8)--(g5) (g3)--(g8) (g7)--(g6) (g1)--(g8) (g4)--(g8) (g2)--(g6);
  \node[accentblue,font=\bfseries\tiny,draw=none,fill=none] at (2.2,3.9) {$G$};

  % tree roots will be at: T1 root=(2.2-6.7, -5.5)=(-4.5,-5.5)
  %                        T2 root=(2.2,      -5.5)
  %                        Tp root=(2.2+6.7,  -5.5)=(8.9,-5.5)
  % arrows - label placed between the three arrows
  \draw[->,mainblue,line width=2pt] (2.2,-2.0) -- (-7.5,-5.0);
  \draw[->,mainblue,line width=2pt] (2.2,-2.0) -- ( 2.2,-5.0);
  \draw[->,mainblue,line width=2pt] (2.2,-2.0) -- (11.9,-5.0);
  \node[mainblue,font=\tiny,draw=none,fill=none] at (2.2,-3.8) {cut-sparsifiers};

  \begin{scope}[shift={(-7.5,-5.5)},accentgreen,line width=1.2pt]
    \node[tn] (t1r)  at ( 0.0, 0.0) {};
    \node[tn] (t1a)  at (-2.0,-1.8) {};
    \node[tn] (t1b)  at ( 2.0,-1.8) {};
    \node[tn] (t1c)  at (-3.0,-3.6) {};
    \node[tn] (t1d)  at (-1.0,-3.6) {};
    \node[tn] (t1e)  at ( 1.0,-3.6) {};
    \node[tn] (t1f)  at ( 3.0,-3.6) {};
    \node[tn] (t1l1) at (-3.5,-5.4) {};
    \node[tn] (t1l2) at (-2.5,-5.4) {};
    \node[tn] (t1l3) at (-1.5,-5.4) {};
    \node[tn] (t1l4) at (-0.5,-5.4) {};
    \node[tn] (t1l5) at ( 0.5,-5.4) {};
    \node[tn] (t1l6) at ( 1.5,-5.4) {};
    \node[tn] (t1l7) at ( 2.5,-5.4) {};
    \node[tn] (t1l8) at ( 3.5,-5.4) {};
    \draw (t1r)--(t1a) (t1r)--(t1b);
    \draw (t1a)--(t1c) (t1a)--(t1d);
    \draw (t1b)--(t1e) (t1b)--(t1f);
    \draw (t1c)--(t1l1) (t1c)--(t1l2);
    \draw (t1d)--(t1l3) (t1d)--(t1l4);
    \draw (t1e)--(t1l5) (t1e)--(t1l6);
    \draw (t1f)--(t1l7) (t1f)--(t1l8);
    \node[font=\bfseries\tiny,draw=none,fill=none,accentgreen] at (0.0,-6.4) {$T_1$};
  \end{scope}

  \begin{scope}[shift={(2.2,-5.5)},accentgreen,line width=1.2pt]
    \node[tn] (t2r)  at ( 0.0, 0.0) {};
    \node[tn] (t2a)  at (-3.5,-1.8) {};
    \node[tn] (t2b)  at ( 0.5,-1.8) {};
    \node[tn] (t2c)  at ( 3.8,-1.8) {};
    \node[tn] (t2d)  at (-0.5,-3.6) {};
    \node[tn] (t2e)  at ( 1.5,-3.6) {};
    \node[tn] (t2l1) at (-4.5,-3.6) {};
    \node[tn] (t2l2) at (-2.5,-3.6) {};
    \node[tn] (t2l3) at (-1.2,-5.4) {};
    \node[tn] (t2l4) at ( 0.2,-5.4) {};
    \node[tn] (t2l5) at ( 0.8,-5.4) {};
    \node[tn] (t2l6) at ( 2.2,-5.4) {};
    \node[tn] (t2l7) at ( 3.0,-3.6) {};
    \node[tn] (t2l8) at ( 4.6,-3.6) {};
    \draw (t2r)--(t2a) (t2r)--(t2b) (t2r)--(t2c);
    \draw (t2a)--(t2l1) (t2a)--(t2l2);
    \draw (t2b)--(t2d) (t2b)--(t2e);
    \draw (t2d)--(t2l3) (t2d)--(t2l4);
    \draw (t2e)--(t2l5) (t2e)--(t2l6);
    \draw (t2c)--(t2l7) (t2c)--(t2l8);
    \node[font=\bfseries\tiny,draw=none,fill=none,accentgreen] at (0.0,-6.4) {$T_2$};
  \end{scope}

  \node[draw=none,fill=none,font=\scriptsize] at (8.1,-7.8) {$\cdots$};

  \begin{scope}[shift={(11.9,-5.5)},accentgreen,line width=1.2pt]
    \node[tn] (tps1) at ( 0.0, 0.0) {};
    \node[tn] (tps2) at ( 0.0,-1.8) {};
    \node[tn] (tps3) at ( 0.0,-3.6) {};
    \node[tn] (tpl1) at (-2.5,-0.9) {};
    \node[tn] (tpl2) at (-1.5,-0.9) {};
    \node[tn] (tpl3) at ( 1.5,-0.9) {};
    \node[tn] (tpl4) at (-1.8,-2.7) {};
    \node[tn] (tpl5) at ( 1.8,-2.7) {};
    \node[tn] (tpl6) at (-2.0,-5.0) {};
    \node[tn] (tpl7) at (-0.5,-5.0) {};
    \node[tn] (tpl8) at ( 1.2,-5.0) {};
    \draw (tps1)--(tps2) (tps2)--(tps3);
    \draw (tps1)--(tpl1) (tps1)--(tpl2) (tps1)--(tpl3);
    \draw (tps2)--(tpl4) (tps2)--(tpl5);
    \draw (tps3)--(tpl6) (tps3)--(tpl7) (tps3)--(tpl8);
    \node[font=\bfseries\tiny,draw=none,fill=none,accentgreen] at (0.0,-6.0) {$T_p$};
  \end{scope}

\end{tikzpicture}}


\vspace{5mm}
\end{tcolorbox}


\end{column}

\begin{column}{0.48\textwidth}

\thmterm{Main Result}{
  \textbf{Theorem.} There is a polynomial-time
  $O(\log n)$-approximation algorithm for minimising $\psi_w(S)$
  on general graphs, and an $O(1)$-approximation on trees.
}


\begin{tcolorbox}[posterbox,title={Main Algorithm (Minimise $\psi_w$)}]
\textbf{Input:} Graph $G=(V,E,c,w)$.\\
\textbf{Output:} Cut $(S,\bar S)$ approximating $\min_S \psi_w(S)$.

\bigskip
\textbf{Step 1 - Candidate $S_1$}
\begin{enumerate}
  \item Build auxiliary graph $G'=(V\cup\{r\},E,c,w')$ by adding a
        root $r$ with $c(r,u)=0$ and $w'(r,u)=w(u,V)$ for every $u\in V$.
  \item Set $K=\tfrac{4}{3}w(V)$.
  \item Run an $\alpha=O\br{\log n}$-approx.\ algorithm for
        \[
          \min_{S\subseteq V'}\phi_{w'}(S)=\frac{c(S,\bar S)}{w'(S,\bar S)}
          \quad\text{s.t.}\quad w'(S,\bar S)\le K.
        \]
  \item Return the corresponding cut $S_1\subseteq V$.
\end{enumerate}

\smallskip\hrule\smallskip

\textbf{Step 2 - Candidate $S_2$:}
\begin{enumerate}
  \item Find a $\beta=O\br{\log n}$-approx.\ $k$-multicut $C$ with $k=w(V)/3$.
  \item Build a complete graph $R$ on components of $G\setminus C$
        with demands $w_R(H_1,H_2)=w_G(H_1,H_2)$.
  \item Run \textbf{local-search max-cut} on $(R,w')$ to obtain a cut $\br{A, B}$:
        start from $(\emptyset,V(R))$, flip a single vertex whenever
        it improves the cut, stop at a local optimum.
  \item Map the cut $\br{A, B}$ of $R$ back to a cut $S_2\subseteq V$ of $G$.
\end{enumerate}

\smallskip\hrule\smallskip

\textbf{Step 3:} Return $\arg\min\bigl\{\psi_w(S_1),\,\psi_w(S_2)\bigr\}$.

\bigskip
\textbf{Key sub-routine:} The constrained sparsest-cut in Step~1 is
solved via \emph{Cut-sparsifier tree decompositions}:
embed $G$ into a convex combination of trees (distortion $O(\log n)$),
solve the problem exactly on each tree, return the best solution.
\end{tcolorbox}


% --- Figure 2: k-multicut + graph R + max-cut ---
\begin{tcolorbox}[posterbox,title={Illustration 4: $k$-Multicut and Max-Cut (Step~2)}]
\small
\textbf{(a)} $k$-multicut partitions $G$ into several components.
\textbf{(b)} Each component becomes a node in a new graph $R$, edge demands equal inter-component demand.
\textbf{(c)} Max-cut on $R$ gives bipartition $\br{A, B}$.

\vspace{3mm}
\centering
\resizebox{1.0\linewidth}{!}{%
\begin{tikzpicture}[>=Stealth,line width=1.2pt,
  iblob/.style={ellipse,draw=mainblue!60,fill=mainblue!10,
          minimum width=9mm,minimum height=6mm,
          line width=0.8pt,font=\tiny,
          preaction={fill=black,opacity=0.18,transform canvas={xshift=0.9pt,yshift=-0.9pt}}},
  rnode/.style={ellipse,draw=mainblue,fill=lightblue!70,
          minimum width=10mm,minimum height=6mm,
          line width=1pt,font=\tiny,
          preaction={fill=black,opacity=0.18,transform canvas={xshift=0.9pt,yshift=-0.9pt}}},
  rA/.style={ellipse,draw=accentblue,fill=accentblue!25,
         minimum width=10mm,minimum height=6mm,
         line width=1pt,font=\tiny,
         preaction={fill=black,opacity=0.18,transform canvas={xshift=0.9pt,yshift=-0.9pt}}},
  rB/.style={ellipse,draw=accentgreen,fill=accentgreen!25,
         minimum width=10mm,minimum height=6mm,
         line width=1pt,font=\tiny,
         preaction={fill=black,opacity=0.18,transform canvas={xshift=0.9pt,yshift=-0.9pt}}}]

  \draw[fill=gray!8,draw=gray!60,line width=2pt,rounded corners=16pt]
    (-0.5,-5.5) rectangle (11.0,6.5);
  \node[draw=none,fill=none,font=\bfseries\tiny,gray!70] at (5.2,7.1) {(a) $G$};

  % Nodes: H5 moved to (0.5,3.5) so H1-H6 edge avoids it
  \node[iblob,fill=accentblue!20,draw=accentblue]       (a1) at (2.0, 5.2) {$H_1$};
  \node[iblob,fill=accentpurple!20,draw=accentpurple!70](a2) at (7.5, 5.2) {$H_2$};
  \node[iblob,fill=accentred!20,draw=accentred]         (a3) at (9.5, 1.5) {$H_3$};
  \node[iblob,fill=accentgreen!20,draw=accentgreen]     (a4) at (5.5, 1.5) {$H_4$};
  \node[iblob,fill=orange!25,draw=orange!80]            (a5) at (0.8,-0.5) {$H_5$};
  \node[iblob,fill=accentblue!15,draw=accentblue!60]    (a6) at (4.5,-3.0) {$H_6$};
  \node[iblob,fill=accentgreen!15,draw=accentgreen!60]  (a7) at (8.0,-2.0) {$H_7$};

  % k-multicut edges with VARYING multiplicity (1, 2, or 3 parallel lines)
  \tikzset{me/.style={accentred,dashed,line width=0.9pt}}
  % 1 edge:
  \draw[me] (a1)--(a2);
  \draw[me] (a2)--(a3);
  \draw[me] (a3)--(a7);
  \draw[me] (a5)--(a6);
  % 2 edges:
  \draw[me] (a1) to[bend left=12]  (a4); \draw[me] (a1) to[bend right=12] (a4);
  \draw[me] (a2) to[bend left=10]  (a4); \draw[me] (a2) to[bend right=10] (a4);
  \draw[me] (a4) to[bend left=12]  (a6); \draw[me] (a4) to[bend right=12] (a6);
  \draw[me] (a6) to[bend left=12]  (a7); \draw[me] (a6) to[bend right=12] (a7);
  % 3 edges:
  \draw[me] (a4) to[bend left=18]  (a5); \draw[me] (a4)--(a5); \draw[me] (a4) to[bend right=18] (a5);
  \node[draw=none,fill=none,font=\tiny,accentred] at (5.0,-5.0)
    {$k$-multicut edges $C$};

  \draw[->,mainblue,line width=3.5pt] (11.2,0.5) -- (13.0,0.5)
    node[midway,above,draw=none,fill=none,font=\tiny,mainblue,align=center]{build\\$R$};

  \node[rnode] (r1) at (16.5, 5.2) {$H_1$};
  \node[rnode] (r2) at (21.5, 5.2) {$H_2$};
  \node[rnode] (r3) at (23.5, 1.5) {$H_3$};
  \node[rnode] (r4) at (19.5, 1.5) {$H_4$};
  \node[rnode] (r5) at (14.8,-0.5) {$H_5$};
  \node[rnode] (r6) at (18.5,-3.0) {$H_6$};
  \node[rnode] (r7) at (22.0,-2.0) {$H_7$};
  % exactly the edges from (a):
  \draw[gray!55,line width=1pt] (r1)--(r2);
  \draw[gray!55,line width=1pt] (r1)--(r4);
  \draw[gray!55,line width=1pt] (r2)--(r3);
  \draw[gray!55,line width=1pt] (r2)--(r4);
  \draw[gray!55,line width=1pt] (r3)--(r7);
  \draw[gray!55,line width=1pt] (r4)--(r5);
  \draw[gray!55,line width=1pt] (r4)--(r6);
  \draw[gray!55,line width=1pt] (r5)--(r6);
  \draw[gray!55,line width=1pt] (r6)--(r7);
  \node[draw=none,fill=none,font=\tiny] at (19.0,3.5){$w_R(H_i,H_j){=}w_G(H_i,H_j)$};
  \node[draw=none,fill=none,font=\bfseries\tiny,gray!70] at (19.0,7.1) {(b) $R$};

  \draw[->,mainblue,line width=3.5pt] (24.8,0.5) -- (26.6,0.5)
    node[midway,above,draw=none,fill=none,font=\tiny,mainblue,align=center]{max-\\cut};

  % Side A (blue): H1,H4,H5  - staggered x positions
  \node[rA] (c1) at (28.5, 5.5) {$H_1$};
  \node[rA] (c4) at (30.0, 2.0) {$H_4$};
  \node[rA] (c5) at (28.5,-1.0) {$H_5$};
  % Side B (green): H2,H3,H6,H7  - staggered
  \node[rB] (c2) at (36.5, 5.5) {$H_2$};
  \node[rB] (c3) at (38.0, 2.5) {$H_3$};
  \node[rB] (c6) at (35.5,-1.5) {$H_6$};
  \node[rB] (c7) at (37.5,-3.0) {$H_7$};
  % cut edges (red) - A={H1,H4,H5}, B={H2,H3,H6,H7}
  % cross edges: a1-a2, a2-a4, a4-a6, a5-a6
  \draw[accentred,line width=2pt] (c1)--(c2);
  \draw[accentred,line width=2pt] (c2)--(c4);
  \draw[accentred,line width=2pt] (c4)--(c6);
  \draw[accentred,line width=2pt] (c5)--(c6);
  % intra-A edges (gray) from (a): a4-a5
  \draw[gray!50,line width=1pt] (c4)--(c5);
  % intra-B edges (gray) from (a): a3-a7, a6-a7
  \draw[gray!50,line width=1pt] (c3)--(c7);
  \draw[gray!50,line width=1pt] (c6)--(c7);
  % non-cut a1-a4 edge is intra-A:
  \draw[gray!50,line width=1pt] (c1)--(c4);
  % non-cut a2-a3 is intra-B:
  \draw[gray!50,line width=1pt] (c2)--(c3);
  \node[accentblue,font=\bfseries\scriptsize,draw=none,fill=none] at (26.8, 2.0) {$A$};
  \node[accentgreen,font=\bfseries\scriptsize,draw=none,fill=none] at (40.0, 2.0) {$B$};
  \draw[dashed,gray!60,line width=2pt] (33.2,-5.5) -- (33.2,7.1);
  \node[draw=none,fill=none,font=\bfseries\tiny,gray!70] at (33.2,7.1){(c) max-cut};

\end{tikzpicture}}

\vspace{2mm}
{\footnotesize
  Red dashed = $k$-multicut edges.  Red solid = max-cut edges of $R$.}
\end{tcolorbox}



%% Reduction map
\begin{tcolorbox}[posterbox,title={Illustration 5: Reduction Map}]
\centering
\resizebox{0.96\linewidth}{!}{%
\begin{tikzpicture}[
  box/.style={draw,rounded corners=3pt,fill=lightblue,
              minimum width=14mm,minimum height=5mm,
              align=center,font=\tiny},
  >=Stealth,line width=0.8pt
]
  % horizontal chain - spread to ~38cm so resizebox scale ~= 1.2x
  \node[box] (csc)   at ( 0.0, 0.0) {Constrained Sparsest\\Cut $\phi_w$};
  \node[box,fill=accentblue!30] (mgc) at (12.0, 0.0) {Min.\ Generalised\\Conductance $\psi_w$};
  \node[box] (gpd)   at (24.0, 0.0) {Graph Partitioning\\with Demands};
  \node[box] (hcd)   at (36.0, 0.0) {Hierarchical Clustering\\with Demands};
  % k-Multicut feeds in from above mgc
  \node[box] (kmcut) at (12.0, 4.0) {$k$-Multicut};

  \draw[->] (csc)   -- (mgc);
  \draw[->] (kmcut) -- (mgc);
  \draw[->] (mgc)   -- (gpd);
  \draw[->] (gpd)   -- (hcd);
  % extend bounding box vertically to match Illustration 1 height
  \path[draw=none] (-1,-1.0) rectangle (40,5.0);
\end{tikzpicture}}
\end{tcolorbox}

\begin{tcolorbox}[appbox,title={Applications}]

\textbf{1.\ Graph Partitioning with Demands.}\\
Find the cheapest edge set $C$, such that every component $H$ of
$G\setminus C$ is small enough, i.e., $w(H)\le w(V)/2$.\\[2pt]
Algorithm: Iteratively apply the $\psi_w$ approximation algorithm, always
removing the smaller side of the returned cut until all components $H$ are of size $w(H)\le 4w(V)/5$.

\smallskip
\thmterm{Corollary}{
  A bicriteria $O(\log n)$-approx. algorithm exists: it
  achieves cost $O(\log n)\cdot c(C^*)$ while guaranteeing
  $w(H)\le \tfrac{4}{5}w(V)$ for every resulting component $H$. A bicriteria $O\br{1}$-approx.  exists for trees.
}

\bigskip
\textbf{2.\ Hierarchical Clustering with Demands.}\\
Find a binary clustering tree $T$ on $V$ minimising
\[
  \sum_{u\in V(T)} w(S_u)\cdot c\bigl(S_u,V\setminus S_u\bigr).
\]
Algorithm: recursively apply graph partitioning with demands,
building $T$ top-down until all clusters become singletons.

\smallskip
\thmterm{Corollary}{
  There is a polynomial-time $O(\log n)$-approximation for
  hierarchical clustering with demands on general graphs, and
  $O(1)$ on trees.
}

\end{tcolorbox}

\end{column}
\end{columns}

\vfill

\begin{tcolorbox}[
  enhanced,colback=titlebg,colframe=titlebg,arc=0pt,boxrule=0pt,
  left=12mm,right=12mm,top=4mm,bottom=4mm
]
  {\color{white}\small
   Szyfelbein \& Dereniowski - Gda\'nsk University of Technology \hfill
   \emph{Graph Partitioning with Demands: Generalized Graph Conductance with Applications to Hierarchical Clustering}}
\end{tcolorbox}

\end{frame}
\end{document}
